\section{Dziennik promptów użytych z modelem}
\label{app:prompty}

Rejestr promptów wysyłanych do modeli podczas tworzenia projektu (kod, dokumentacja, raport).
Wymóg przedmiotu: dokumentować najważniejsze prompty, iteracje, modele i ustawienia. Wpisy
dodajemy chronologicznie na bieżąco (zasady w \texttt{CLAUDE.md}). Treści wygenerowane przez
AI są oznaczone w raporcie, a wyniki weryfikuje zespół.

Nie obejmuje promptów systemowych samego asystenta RAG: te są w \texttt{prompts/} i opisane w
sekcji o strategii promptowania.

% ---------------------------------------------------------------- wpisy

\subsection*{Etap wyboru tematu}

Temat wybrano w rozmowie z Gemini (burza mózgów, 10 kandydatów, wybór tematu~1). Wpisy 1 i~2
dokumentują ten etap. Dokładne wersje modelu i ustawienia nie są znane (czat webowy).

\promptentry{1}{nieustalona (eksport rozmowy: 2026-10-28)}{Gemini (wersja nieustalona), czat webowy}{nieustalone}
  {Burza mózgów: 10 tematów dla projektu RAG, roadmapy i stacki}
\begin{lstlisting}
Przeanalizuj mi ten projekt robie z ziomalem temat 12 czyli poprostu RAG najbardziej
ciekawszy, zrob brainstorming wypisz mi z 10 tematow pod co tego raga mozna zrobic jaka
tematyke uzyc aby dobrze pokazac karmienie llma ragiem takie z ragiem i bez raga jak to
wyglada albo jakis chatbot na jakiejs stronie albo cos w tym stylu, odpal research exa
mcp i zrob mi brainstorming odnosnie tematow i wypisz do kazdego krociutka road mapke jak
do tego podejsc i stack oczywiscie. Prowadzacy polecal jakas biblioteke pythonowska do
zrobienia tego raga ale nie pamietam [...]
\end{lstlisting}
\promptresult{Lista 10 tematów z krótką roadmapą i stackiem. Wskazanie Chainlit, Streamlit i
Gradio jako bibliotek UI. Zestawienie wymagań zaliczeniowych.}

\promptentry{2}{nieustalona (eksport rozmowy: 2026-10-28)}{Gemini (wersja nieustalona), czat webowy}{nieustalone}
  {Wybór tematu~1 (sylabusy), doprecyzowanie UI, darmowych modeli i scenariuszy odmowy}
\begin{lstlisting}
1. Uczelniany Asystent Sylabusow i Zasad Zaliczen [opis tematu z poprzedniej odpowiedzi]
Dobra ewidentnie bierzemy to, rozwin ten temat powiedz stack jak bedzie UI wygladac czy
to sam chatbot poprostu zrobimy czy jakis front w react co myslisz wg tego projektu i
wytycznych daj znac. Dodatkowo tak polecal on gradio ale daj znac ktory jest najlepszy
chainlit czy streamlit [...] czy z openroutera da sie wziac api darmowego modelu zeby
potestowac tego raga [...] Dodatkowo jak zlapac tego raga na tym ze nie bedzie chcial
odpowiedziec na wiadomosc w danym temacie [...]
\end{lstlisting}
\promptresult{Temat~1 wybrany. Rekomendacja gotowej biblioteki Pythona zamiast frontu w React,
wstępny faworyt Chainlit, darmowe modele z OpenRoutera, lokalne embeddingi, trzy scenariusze
testowe odmowy. Rekomendacje do weryfikacji w etapie wyboru stacku.}

\subsection*{Etap inicjalizacji projektu}

\promptentry{3}{2026-10-28}{Claude Sonnet 5.5 (\texttt{claude-sonnet-5-5}), Claude Code}{medium}
  {Analiza pomysłu projektu, inicjalizacja repozytorium, struktura plików, szkielet raportu LaTeX}
\begin{lstlisting}
Przenalizuj mi idea calego projektu jaki tu rozpisalem (PDF: Tematy i roadmapy
projektow RAG + zrzuty wymagan przedmiotu), dodatkowo wez moja rozmowe z gemini
w ktorym dyskutowalem na temat projektu, zacznij od inicjalizacji repo daj maina
zrob go narazie publicznie, stworz README oraz przygotuj jakas sensowna strukture
plikow, cos w stylu folder raport gdzie bedzie miejsce do generowania latexa z
raportem promptow jakie uzyte z modelem; do claude.md dopisz aby w tym pliku
latex odrazu spisywal prompty uzyte z modelem i effortem, ta rozmowe z gemini
tez troche wez i dopisz juz na start ten proces jak klepnelismy temat
1. Uczelniany Asystent Sylabusow i Zasad Zaliczen [opis tematu z rozmowy: bez RAG
vs z RAG, roadmapa, stack Python/LangChain/PyPDFLoader/ChromaDB/Streamlit-Chainlit]
i opisz mi to i zaczniemy sobie od takiej wstepnej inicjalizacji projektu, a
stack w nastepnym etapie sie zajmiemy po analizie gradio/streamlit, no
przejrzymy.
\end{lstlisting}
\promptresult{Repozytorium z gałęzią \texttt{main}, README, \texttt{docs/} (wymagania, opis
projektu, plan), szkielet raportu w \texttt{raport/}, \texttt{CLAUDE.md} z regułą dziennika
promptów. Wybór stacku odłożony do następnego etapu.}
